\section{Обзор существующих решений}
\label{sec:Chapter2} \index{Chapter2}

На данный момент большинство виртуальных машин обладают верификаторами, в качестве самых популярных 
примеров могут послужить JVM - виртуальная машина для исполнения байт-кода таких языков как: java, kotlin;
.Net CLR - виртуальная машина для исполнения байт-кода С\#, F\#. Так же существую верификаторы, которые
не поставляются вместе с платформой, а являются сторонними библотеками, примерами могут послужить 
BCEL Verifier для Java байт-кода и PEVerify, который проверяет на корректность IL-кода(С\#). 
Рассмотрим некоторые из существующих решений и виды выполняемых проверок более подробно.

\section{Верификация байт-кода динамически типизированных яп:}
В языках с динамической типизацией типы определяются во время выполнения программы: переменная может 
изменять свой тип, а прототипы могут модифицироваться динамически. К тому же во многих языках 
есть конструкции, принимающие строки 
кода и выполняющие их уже в рантайме. В свзязи с этим список возможных проверок ограничен.

\begin{enumerate}

    \item \textbf{Проверки формата байткода:} 
        
    Целью данных проверок является тестирование на валидность инструкций, корректность их формата и 
    правильность индексов, соответствующих методам, полям и типам.

    \item \textbf{Проверки графа управления} 

    Во время выполнения данного типа проверок, программа анализируется 
    на корректность инструкций условного перехода: целевым адресом перехода не выходит за пределы 
    текущего метода, не является серединой try-блока серединой другой инструкции. Проверяется отсутствие циклов 
    с некорректным состоянием регистров. Контролируется, что все пути выполнения заканчиваются 
    инструкциями обозначающими $return$ или $throw$

    \item \textbf{Сохранение инварианта стека}
    
    Размер стека должен быть неотрицательным числом, при этом его глубина не должна привышать некоторого
    заданного значения, в точках схождения глубина стекак должна быть одинаковой.

    \item \textbf{Проверка Инлайн кэшей:}
    
    Inline Cache — это механизм ускорения динамического доступа в JS. Но для безопасного их использования
    нам надо гарантировать, что при обращении к конкретному слоту он будет существовать и будет 
    инициализирован, а так же, что данный слот будет соответствовать типу текущей инструции.

\end{enumerate}

Важно отметить, что как правило движки динамически типизированных яп в себе не имеют такого объекта как 
верификатор, а перечисленные проверки выполняются непосредственно на этапе выполнения.

\section{Верификация байт-кода статически типизированных яп:}
В статически типизированных языках типы переменных, аргументов и возвращаемых значений известны на 
этапе компиляции. Это позволяет компилятору генерировать байткод с встроенными гарантиями типов.
Но байткод может быть сгенерирован не только компилятором, поэтому верификатор остается неотъемлимой частью
вертуальной машины. Однако теперь возможно выполнения ряда новых проверок, в добавок к тем, что уже были 
перечислены ранее:

\begin{enumerate}
    \item \textbf{Проверки типовой безопастности регистров:}
    
    Байткод инспектируется на предмет того, что тип каждого регистра является примитивным или 
    ссылочным. Верифицируется, что любой регист используется после своей инициализации.

    \item \textbf{Проверки инструкций вызова:}
    
    Инструкции вызова проходят проверку на корректность объекта ресивера(в случае нестатического метода).
    На количество аргументов, на типы данных аргументов и тип возвращаемого значения. Верифицируется
    возможность доступа к полю или методу класса.

    \item \textbf{Проверка доступа обращения:}
    
    Именно данная верификация является одной из самых важных для безопасности. Как известно, 
    существует 3 базовых модификатора доступа: $публичные$, $защищенные$ и $приватные$, которые
    применимы, как к элементам класса, так и классам и целым модулям. Контроль того, что их семантика
    не нарушается на уровне байткода необходима.
\end{enumerate}

Рассмотрим пару примеров промышленных решений.
\begin{enumerate}

    \item \textbf{ART}

    \text{ART - Android Runtime}
    В контексте данной работы для нас наибольший интерес представляют верификаторы тех платформ, которые 
    предназначены для исполнения байт-кода на мобильных устройствах. Именно таким является верификатор ART(Android
    runtime). ART исполняет DEX-байткод, источником которого могут быть: сторонние приложения, 
    сгенерированного кода, потенциально повреждённых или злонамеренных источников. 
    \item \textbf{LLVM} 

    \textit{LLVM - Low level virtual machine}
    LLVM-верификация — это обладает несколько иной философия, чем ART: она не про безопасность 
    пользователя, а про корректность как ПП(здесь и далее ПП - промежуточное представление) компилятора. ПП 
    генерируется доверенным фронтендом, но оптимизации могут его сломать. Можно выделить 4 уровня верификации:
    уровень модуля, функций, базовых блоков и инструкций. Проверяются как локальные, так глобальные инварианты.
\end{enumerate}
